\section{Introdução}

Streaming ao vivo combina requisitos de baixa latência, estabilidade operacional e eficiência no uso de recursos. Em cenários acadêmicos e de produção, pequenas mudanças no servidor de ingestão ou no transcodificador podem alterar o comportamento observável pelo espectador, principalmente quando há múltiplos perfis de saída e vários consumidores simultâneos. Este projeto parte dessa premissa para comparar, em condições controladas, duas opções de servidor RTMP e dois transcodificadores amplamente adotados na prática.

A contribuição inicial do trabalho não está apenas na implementação da plataforma, mas também na organização dos artefatos para escrita científica. Os resultados brutos, os agregados estatísticos, os logs e o relatório consolidado foram estruturados para que cada seção do TCC possa ser derivada do mesmo conjunto de evidências. Isso reduz retrabalho e melhora a rastreabilidade entre experimento, análise e redação.

Do ponto de vista técnico, a versão atual da solução trabalha com transcodificação em software. FFmpeg e GStreamer aparecem como opções alternativas para o mesmo fluxo de entrega HLS, mas sem uso de GPU no pipeline. Esse detalhe é importante porque a interpretação de CPU e memória depende do host, enquanto a comparação entre cenários deve ocorrer dentro do mesmo ambiente documentado.

\subsection{Objetivo geral}

Avaliar e comparar o comportamento de servidores RTMP e transcodificadores em uma plataforma de streaming ao vivo, observando latência, consumo de recursos, erros de entrega e estabilidade operacional.

\subsection{Objetivos específicos}

\begin{itemize}[leftmargin=1.2cm]
  \item Documentar a arquitetura de software e infraestrutura usada no experimento.
  \item Organizar os dados de benchmark em formato apropriado para escrita acadêmica.
  \item Comparar Nginx-RTMP e SRS com FFmpeg e GStreamer sob cargas distintas.
  \item Consolidar os resultados recentes em tabelas que possam ser reaproveitadas no texto final.
\end{itemize}